%==============================================================================
%  Relazione di progetto -- Miglioramento di TabPFN-2.5 su classificazione
%  binaria: adattamento interno (LoRA) e pipeline esterna.
%  Compilare con: pdflatex relazione_completa.tex  (due volte per i riferimenti)
%  Richiede nella stessa cartella: capitolo_pipeline_esterna.tex
%==============================================================================
\documentclass[11pt,a4paper]{article}

\usepackage[utf8]{inputenc}
\usepackage[T1]{fontenc}
\usepackage[italian]{babel}
\usepackage{lmodern}
\usepackage{microtype}
\usepackage[margin=2.5cm]{geometry}
\usepackage{amsmath,amssymb}
\usepackage{booktabs}
\usepackage{multirow}
\usepackage{graphicx}
\usepackage{xcolor}
\usepackage{listings}
\usepackage{caption}
\usepackage{enumitem}
\usepackage{pgfplots}
\pgfplotsset{compat=1.17}
\usepackage[hidelinks]{hyperref}

% --- Stile per i blocchi di codice ------------------------------------------
\definecolor{codegray}{gray}{0.95}
\definecolor{codekw}{rgb}{0.13,0.29,0.53}
\definecolor{codecmt}{rgb}{0.0,0.5,0.0}
\definecolor{codestr}{rgb}{0.58,0.0,0.0}
\lstset{
  basicstyle=\ttfamily\footnotesize,
  backgroundcolor=\color{codegray},
  keywordstyle=\color{codekw}\bfseries,
  commentstyle=\color{codecmt}\itshape,
  stringstyle=\color{codestr},
  numbers=left,
  numberstyle=\tiny\color{gray},
  frame=single,
  framesep=4pt,
  breaklines=true,
  showstringspaces=false,
  language=Python,
  captionpos=b,
}

\title{\textbf{Migliorare TabPFN-2.5 su classificazione binaria}\\
\large Adattamento interno tramite LoRA e pipeline esterna di
ensembling, calibrazione e soglia}
\author{Progetto di Machine Learning}
\date{\today}

\begin{document}
\maketitle

\begin{abstract}
\noindent
Questo lavoro studia se e come sia possibile migliorare le prestazioni del
\emph{tabular foundation model} TabPFN-2.5 su task di classificazione binaria,
sotto il vincolo di \emph{non} riaddestrarne i pesi principali. Vengono
esplorate due direzioni complementari. La prima, \emph{interna}, è un
fine-tuning leggero basato su \emph{Low-Rank Adaptation} (LoRA): si descrive
l'implementazione completa (iniezione e addestramento degli adapter) e si
conducono tre esperimenti su dataset del dominio finanziario/creditizio --
generalizzazione \emph{in-domain}, ablation sulla capacità degli adapter e
ablation sulla dimensione dei dati. La seconda, \emph{esterna},
lascia i pesi congelati e agisce sull'output del modello: ensembling
(self-ensemble ed eterogeneo), calibrazione delle probabilità (Platt e
isotonica) e ottimizzazione della soglia decisionale per l'F1 macro, valutate
con ripetizioni multiple e split a tre vie anti-leakage. I risultati convergono verso
una tesi unica: TabPFN-2.5 è vicino all'ottimo non solo nella capacità
discriminativa, ma anche nella calibrazione e nella varianza delle stime. Né il
LoRA (anche in-domain, con $4\times$ capacità o dati fino a $80\times$), né
l'ensembling, né la ricalibrazione producono miglioramenti robusti. L'unico
intervento con effetto
reale è l'ottimizzazione della soglia, che migliora l'F1 macro di $4$--$5$ punti
sui dataset sbilanciati: un guadagno che non tocca né i pesi né le probabilità,
ma soltanto la regola di decisione. Si tratta di un risultato negativo per le
tecniche di adattamento, ma metodologicamente solido e con un'indicazione
pratica netta.
\end{abstract}

\tableofcontents
\bigskip

%==============================================================================
\section{Introduzione}
%==============================================================================

\subsection{Contesto: TabPFN-2.5}
TabPFN è un \emph{tabular foundation model} che affronta la classificazione su
dati tabulari mediante \emph{in-context learning}: invece di essere addestrato
su uno specifico dataset, il modello riceve in input l'intero training set come
``contesto'' e produce le predizioni per le righe di test in un singolo forward
pass, senza ottimizzazione dei pesi a tempo di utilizzo. I pesi del modello
sono pre-addestrati una volta per tutte su milioni di dataset sintetici. In
questo progetto si utilizza la versione \textbf{2.5} (non la v3, che è il
default del pacchetto), che supporta dataset fino a circa $50\,000$ campioni.

\subsection{Obiettivo e struttura della relazione}
L'obiettivo è migliorare le prestazioni di TabPFN-2.5 su classificazione binaria
senza riaddestrare i pesi principali, esplorando due famiglie di tecniche:
\begin{itemize}[nosep]
  \item un \textbf{adattamento interno} leggero tramite LoRA fine-tuning
  (Sez.~\ref{sec:lora});
  \item una \textbf{pipeline esterna} che agisce sull'output -- ensembling,
  calibrazione e ottimizzazione della soglia (Sez.~\ref{sec:pipeline}).
\end{itemize}
La Sezione~\ref{sec:metodi} descrive i materiali e i metodi comuni (dati,
metriche, architettura e implementazione). I due blocchi sperimentali sono
volutamente speculari: il primo modifica il modello, il secondo lo lascia
intatto e ne post-processa le uscite. La Sezione~\ref{sec:conclusioni} ne trae
le conclusioni complessive.

\subsection{Low-Rank Adaptation (LoRA)}
LoRA è una tecnica di \emph{parameter-efficient fine-tuning}: dato un layer
lineare con matrice di pesi $W \in \mathbb{R}^{d_{out}\times d_{in}}$, invece di
aggiornare $W$ si congela e si aggiunge un termine a basso rango
\begin{equation}
  y = Wx + \frac{\alpha}{r}\, B A\, x,
  \qquad A \in \mathbb{R}^{r\times d_{in}},\ B \in \mathbb{R}^{d_{out}\times r},
  \label{eq:lora}
\end{equation}
dove $r \ll \min(d_{in},d_{out})$ è il \emph{rango} e $\alpha$ un fattore di
scala. Si addestrano soltanto le matrici $A$ e $B$. All'inizializzazione si pone
$B = 0$, cosicché il termine aggiuntivo è nullo e il modello parte
numericamente identico a quello pre-addestrato; l'adattamento emerge solo
durante l'addestramento. Il numero di parametri allenabili passa così da
$d_{out}\,d_{in}$ a $r\,(d_{in}+d_{out})$, tipicamente una piccola frazione del
totale.

%==============================================================================
\section{Materiali e metodi}
\label{sec:metodi}
%==============================================================================

\subsection{Dataset e preprocessing}
Il fine-tuning LoRA impiega dataset del \textbf{dominio finanziario/creditizio}
(Tabella~\ref{tab:dataset}): cinque dataset per l'addestramento e quattro per la
valutazione, mai usati in training. La scelta di restare all'interno di un unico
dominio consente una valutazione \emph{in-domain} (finanziario~$\to$~finanziario),
adatta a rispondere alla domanda ``il LoRA riesce a specializzare TabPFN su uno
specifico settore?''. L'ablation sulla dimensione dei dati (Sez.~\ref{sec:datasize})
usa inoltre il dataset \texttt{gmsc} (\emph{Give-Me-Some-Credit}). La pipeline
esterna (Sez.~\ref{sec:pipeline}) adotta un proprio benchmark di sei dataset,
descritto nella Sez.~\ref{sec:pipeline-metodo}.

\begin{table}[h]
\centering
\caption{Dataset finanziari utilizzati per il LoRA (OpenML). I dataset di
valutazione non sono mai stati impiegati per l'addestramento degli adapter.}
\label{tab:dataset}
\small
\begin{tabular}{llrr}
\toprule
Ruolo & Dataset (OpenML ID) & Campioni & Feature \\
\midrule
Fine-tuning & \texttt{bank\_marketing} (1461)      & 45211 & 16 \\
Fine-tuning & \texttt{default\_credit} (42477)     & 30000 & 23 \\
Fine-tuning & \texttt{polish\_bankruptcy\_2} (42984) & 10173 & 64 \\
Fine-tuning & \texttt{polish\_bankruptcy\_3} (42985) & 10503 & 64 \\
Fine-tuning & \texttt{polish\_bankruptcy\_4} (42986) & 9792  & 64 \\
\midrule
Valutazione & \texttt{polish\_bankruptcy\_1} (42880) & 7027  & 64 \\
Valutazione & \texttt{polish\_bankruptcy\_5} (42987) & 5910  & 64 \\
Valutazione & \texttt{australian\_credit} (40981)  & 690   & 14 \\
Valutazione & \texttt{credit\_g} (31)              & 1000  & 20 \\
\midrule
Ablation dati & \texttt{gmsc} (45577)              & 150000 & 10 \\
\bottomrule
\end{tabular}
\end{table}

\paragraph{Verifica degli identificativi.} Alcuni ID OpenML forniti nelle
specifiche iniziali si sono rivelati errati e sono stati corretti dopo verifica
diretta (al caricamento ogni dataset stampa nome, dimensioni e distribuzione del
target, come controllo):
\begin{itemize}[nosep]
  \item i dataset \texttt{polish\_bankruptcy\_1--5}: gli ID 40474--40478 indicati
  corrispondono in realtà a dataset \texttt{thyroid} a \emph{cinque} classi (non
  binari); gli ID corretti dei \texttt{polish-bankruptcy-Nyear} sono
  \textbf{42880}, \textbf{42984}, \textbf{42985}, \textbf{42986}, \textbf{42987};
  \item \texttt{gmsc}: l'ID 44089 conta solo $16\,714$ righe (insufficienti per un
  training fino a $40\,000$); l'ID corretto è \textbf{45577}
  (\texttt{Give-Me-Some-Credit}, $150\,000$ righe).
\end{itemize}
Si noti che i dataset \texttt{polish\_bankruptcy} sono fortemente sbilanciati
(classe positiva $4$--$7\%$), il che rende particolarmente informative l'F1 macro
e l'ottimizzazione della soglia.

\paragraph{Pipeline di caricamento.} È stato implementato un modulo
(\texttt{utils/data\_loader.py}) che, per ciascun dataset: (1)~scarica i dati da
OpenML con caching locale; (2)~codifica le feature categoriche con
\texttt{OrdinalEncoder}, distinguendo numerico e categorico in modo robusto al
\emph{backend} dei tipi (in \texttt{pandas}~3.0 le stringhe hanno dtype
\texttt{str} Arrow, non \texttt{object}); (3)~imputa i valori mancanti con la
mediana (\texttt{SimpleImputer}); (4)~codifica il target in formato binario
$\{0,1\}$ con \texttt{LabelEncoder}, rimuovendo spazi/tabulazioni accidentali;
(5)~scarta colonne identificative (es. \texttt{id}) che introdurrebbero
\emph{data leakage}. La suddivisione è sempre stratificata.

\subsection{Metriche di valutazione}
Sono state implementate (modulo \texttt{evaluation/metrics.py}) quattro metriche
per la classificazione binaria, due di discriminazione e due di calibrazione:
\begin{itemize}[nosep]
  \item \textbf{AUC-ROC}: area sotto la curva ROC, capacità di discriminazione;
  \item \textbf{F1 macro}: media non pesata dell'F1 per classe, adatta a dataset
  sbilanciati;
  \item \textbf{Brier score}: $\frac{1}{N}\sum_i (p_i - y_i)^2$, errore quadratico
  medio delle probabilità (calibrazione);
  \item \textbf{Expected Calibration Error (ECE)}: suddividendo le probabilità in
  $M=10$ bin di uguale ampiezza,
  \begin{equation}
    \mathrm{ECE} = \sum_{b=1}^{M} \frac{|B_b|}{N}\,
    \bigl|\,\overline{p}(B_b) - \overline{y}(B_b)\,\bigr|,
  \end{equation}
  dove $\overline{p}(B_b)$ e $\overline{y}(B_b)$ sono la confidenza media e
  l'accuratezza media nel bin $b$.
\end{itemize}

\subsection{Architettura di TabPFN-2.5 e punti di iniezione}
\label{sec:arch}
L'architettura (file \texttt{architectures/tabpfn\_v2\_5.py}) è composta da una
pila di $L=24$ blocchi identici. Ogni blocco contiene due moduli di attenzione e
un MLP:
\begin{itemize}[nosep]
  \item \texttt{per\_sample\_attention\_between\_features}: attenzione tra le
  feature di una stessa riga;
  \item \texttt{per\_column\_attention\_between\_cells}: attenzione tra le righe
  di una stessa colonna (con \emph{multi-query attention} per le righe di test).
\end{itemize}
Entrambe ereditano da una classe \texttt{Attention} che definisce quattro
proiezioni lineari \emph{senza bias}: \texttt{q\_projection},
\texttt{k\_projection}, \texttt{v\_projection}, \texttt{out\_projection}, tutte
$\mathbb{R}^{192\times192}$ (dimensione di embedding $d=192$, $3$ teste,
\texttt{head\_dim}$=64$). Il modello completo conta circa $11{,}0$ milioni di
parametri.

I \textbf{punti di iniezione naturali} per LoRA sono dunque le proiezioni
dell'attenzione. Adottando per default il classico schema ``QV'' (solo
\texttt{q\_projection} e \texttt{v\_projection}) si ottengono
\[
  L \times 2~\text{(attenzioni)} \times 2~\text{(proiezioni)} = 24\times2\times2 = 96
  \quad\text{adapter}.
\]

\subsection{Implementazione del modulo LoRA}
Il modulo \texttt{models/lora.py} implementa l'Eq.~\eqref{eq:lora} come
\emph{wrapper} attorno a un \texttt{nn.Linear} esistente, indipendentemente da
TabPFN (e quindi testabile in isolamento). Il nucleo è riportato nel
Listing~\ref{lst:lora}.

\begin{lstlisting}[caption={Nucleo del wrapper \texttt{LoRALinear} (semplificato).},label=lst:lora]
class LoRALinear(nn.Module):
    def __init__(self, base_layer, r=8, alpha=16.0, dropout=0.0):
        super().__init__()
        self.base = base_layer                 # congelato
        for p in self.base.parameters():
            p.requires_grad = False
        self.scaling = alpha / r
        self.lora_A = nn.Parameter(torch.empty(r, base_layer.in_features))
        self.lora_B = nn.Parameter(torch.empty(base_layer.out_features, r))
        nn.init.kaiming_uniform_(self.lora_A, a=math.sqrt(5))
        nn.init.zeros_(self.lora_B)            # B=0 -> parte come il modello base

    def forward(self, x):
        lora = (self.lora_dropout(x) @ self.lora_A.t()) @ self.lora_B.t()
        return self.base(x) + self.scaling * lora
\end{lstlisting}

Sono inoltre fornite funzioni per: iniettare gli adapter individuando per nome i
\texttt{nn.Linear} bersaglio (\texttt{inject\_lora\_adapters}); congelare tutti i
parametri tranne quelli LoRA (\texttt{mark\_only\_lora\_as\_trainable});
salvare/caricare i soli pesi degli adapter (pochi MB); e fondere gli adapter nei
pesi base per l'inferenza (\texttt{merge}). Il modulo è stato validato con test
unitari su CPU che verificano, fra l'altro, l'invarianza dell'output
all'inizializzazione ($\lVert y_{\text{base}}-y_{\text{LoRA}}\rVert_\infty = 0$)
e la corretta riproducibilità dopo salvataggio e ricarica degli adapter.

\subsection{Procedura di addestramento LoRA}
L'addestramento (modulo \texttt{models/tabpfn\_lora.py}) adotta un \emph{loop
custom e trasparente}, che riutilizza le primitive di preprocessing di TabPFN ma
controlla esplicitamente quali parametri vengono ottimizzati. Per ogni
\emph{meta-batch} (corrispondente a un dataset, suddiviso in contesto e query):
\begin{enumerate}[nosep]
  \item si imposta il contesto in-context (\texttt{fit\_from\_preprocessed});
  \item si calcolano i \emph{logits} grezzi sulle query
  (\texttt{forward(\dots, return\_raw\_logits=True)});
  \item si calcola la cross-entropy tra logits e label delle query;
  \item si esegue \texttt{backward} e un passo di ottimizzazione,
  \emph{aggiornando soltanto i parametri LoRA}.
\end{enumerate}
L'ottimizzatore è AdamW con \emph{learning rate} $10^{-4}$ e \emph{gradient
clipping}. Per evitare l'esaurimento della memoria GPU su contesti ampi si
abilita l'\emph{activation checkpointing} (ricalcolo delle attivazioni in fase di
backward anziché memorizzazione).

\paragraph{Early stopping.} Una porzione del training set (separata dal
validation usato per il confronto finale, per evitare \emph{leakage}) viene
riservata alla validazione interna: a ogni epoca si misura l'AUC, si conservano i
pesi LoRA migliori e ci si arresta dopo un numero prefissato di epoche senza
miglioramento, ripristinando il miglior checkpoint. Questo accorgimento si è
rivelato essenziale (cfr.\ Sez.~\ref{sec:datasize}).

%==============================================================================
\section{Adattamento interno: LoRA fine-tuning}
\label{sec:lora}
%==============================================================================

Gli esperimenti di questa sezione sono stati eseguiti su GPU (Google Colab,
NVIDIA T4). Baseline e modello LoRA condividono lo stesso numero di estimatori in
inferenza per garantire un confronto equo.

\subsection{Esperimento 1: generalizzazione in-domain}
\label{sec:exp5}
Un \emph{unico} adapter LoRA ($r=8$, $\alpha=16$, schema QV; $294{,}912$
parametri allenabili pari al $2{,}68\%$ del totale) è stato addestrato
congiuntamente sui cinque dataset finanziari di fine-tuning e poi valutato sui
quattro dataset finanziari di test mai visti, per misurare la generalizzazione
\emph{in-domain} dell'adattamento. La migliore AUC media di validazione
($0{,}911$) è stata raggiunta già alla seconda epoca, per poi decrescere: un
primo segnale che l'addestramento allontana l'adapter dal buon punto di partenza.
I risultati sui dataset di test sono riportati in Tabella~\ref{tab:exp5}.

\begin{table}[h]
\centering
\caption{Esperimento 1 -- TabPFN base vs.\ TabPFN+LoRA sui dataset di
valutazione in-domain. In \textbf{grassetto} il valore migliore per metrica e
dataset. Frecce: $\uparrow$ meglio se alto, $\downarrow$ meglio se basso.}
\label{tab:exp5}
\small
\begin{tabular}{llcccc}
\toprule
Dataset & Modello & AUC-ROC$\uparrow$ & F1$\uparrow$ & Brier$\downarrow$ & ECE$\downarrow$ \\
\midrule
\multirow{2}{*}{\texttt{australian\_credit}}
 & base & 0.9361 & \textbf{0.8261} & \textbf{0.1080} & \textbf{0.0686} \\
 & LoRA & \textbf{0.9373} & 0.8188 & \textbf{0.1080} & 0.0709 \\
\addlinespace
\multirow{2}{*}{\texttt{credit\_g}}
 & base & 0.7750 & \textbf{0.6796} & 0.1677 & 0.0399 \\
 & LoRA & \textbf{0.7758} & \textbf{0.6796} & \textbf{0.1675} & \textbf{0.0372} \\
\addlinespace
\multirow{2}{*}{\texttt{polish\_bankruptcy\_1}}
 & base & \textbf{0.9932} & \textbf{0.9055} & \textbf{0.0114} & 0.0193 \\
 & LoRA & 0.9930 & \textbf{0.9055} & 0.0115 & \textbf{0.0191} \\
\addlinespace
\multirow{2}{*}{\texttt{polish\_bankruptcy\_5}}
 & base & \textbf{0.9739} & 0.8518 & \textbf{0.0249} & \textbf{0.0161} \\
 & LoRA & 0.9734 & \textbf{0.8547} & 0.0250 & 0.0166 \\
\bottomrule
\end{tabular}
\end{table}

Le differenze sono tutte dell'ordine della terza--quarta cifra decimale, ovvero
entro il rumore statistico: il LoRA \emph{non} migliora né peggiora in modo
significativo alcuna metrica. Il risultato è più forte di quanto si otterrebbe
in un setting cross-domain, perché qui addestramento e valutazione condividono
lo stesso dominio: nemmeno \emph{in-domain} l'adattamento riesce a specializzare
TabPFN. Si noti che i dataset \texttt{polish\_bankruptcy}, molto sbilanciati,
hanno AUC molto alta ($\approx 0{,}97$--$0{,}99$) già per il modello base.

\subsection{Esperimento 2: ablation sulla capacità}
\label{sec:capacity}
Per verificare se l'assenza di miglioramento dipenda da una capacità
insufficiente degli adapter, l'Esperimento~1 è stato ripetuto con tre
configurazioni crescenti (Tabella~\ref{tab:capacity}).

\begin{table}[h]
\centering
\caption{Esperimento 2 -- configurazioni LoRA a capacità crescente. La migliore
AUC di validazione resta sostanzialmente costante.}
\label{tab:capacity}
\small
\begin{tabular}{lrrc}
\toprule
Config. & Target & Param.\ allenabili & Best val AUC \\
\midrule
\texttt{r8\_qv}    & q, v       & $294{,}912$\ ($2{,}7\%$)   & 0.9108 \\
\texttt{r16\_qv}   & q, v       & $589{,}824$\ ($5{,}2\%$)   & 0.9105 \\
\texttt{r16\_qkvo} & q, k, v, o & $1{,}179{,}648$\ ($9{,}9\%$) & 0.9100 \\
\bottomrule
\end{tabular}
\end{table}

Quadruplicando i parametri allenabili (dal $2{,}7\%$ al $9{,}9\%$) la migliore
AUC di validazione resta sostanzialmente invariata ($\approx 0{,}910$, con anzi
un lievissimo calo) e le metriche sui dataset di test sono indistinguibili. La
configurazione più capiente (\texttt{r16\_qkvo}) peggiora l'F1 su
\texttt{polish\_bankruptcy\_1} ($0{,}9055 \to 0{,}8921$), segno di un principio di
\emph{overfitting}: più capacità non solo non aiuta, ma comincia a danneggiare.
Si conclude che \textbf{la capacità degli adapter non è il fattore limitante}.

\subsection{Esperimento 3: ablation sulla dimensione dei dati}
\label{sec:datasize}
L'ultima ipotesi alternativa è che il LoRA non migliori per \emph{scarsità di
dati} di addestramento. Per testarla in-domain si usa \texttt{gmsc}
(\emph{Give-Me-Some-Credit}, OpenML~45577, $150\,000$ righe). Si ritaglia un
\textbf{test set fisso} di $8\,000$ righe (mai modificato) e si allena l'adapter
($r=8$, QV) su training set di dimensione crescente campionati stratificati dal
resto: $500$, $2\,000$, $10\,000$ e $40\,000$ righe. Ogni configurazione rispetta
il limite di TabPFN ($40\,000 + 8\,000 < 50\,000$). Il confronto base vs.\ LoRA
sul medesimo test è riportato in Tabella~\ref{tab:datasize} e
Figura~\ref{fig:datasize}.

\begin{table}[h]
\centering
\caption{Esperimento 3 -- TabPFN base vs.\ LoRA su \texttt{gmsc} al crescere del
training set (test fisso di $8\,000$ righe).}
\label{tab:datasize}
\small
\begin{tabular}{llcccc}
\toprule
Training & Modello & AUC-ROC$\uparrow$ & F1$\uparrow$ & Brier$\downarrow$ & ECE$\downarrow$ \\
\midrule
\multirow{2}{*}{$500$}
 & base & \textbf{0.8485} & \textbf{0.5362} & \textbf{0.0520} & 0.0119 \\
 & LoRA & 0.8481 & 0.5330 & 0.0521 & \textbf{0.0116} \\
\addlinespace
\multirow{2}{*}{$2\,000$}
 & base & \textbf{0.8565} & \textbf{0.5623} & \textbf{0.0505} & \textbf{0.0099} \\
 & LoRA & 0.8553 & 0.5609 & 0.0506 & 0.0100 \\
\addlinespace
\multirow{2}{*}{$10\,000$}
 & base & \textbf{0.8658} & 0.5786 & \textbf{0.0495} & \textbf{0.0090} \\
 & LoRA & 0.8657 & \textbf{0.5788} & \textbf{0.0495} & 0.0091 \\
\addlinespace
\multirow{2}{*}{$40\,000$}
 & base & \textbf{0.8690} & \textbf{0.6268} & \textbf{0.0491} & \textbf{0.0074} \\
 & LoRA & 0.8674 & 0.6156 & 0.0494 & 0.0163 \\
\bottomrule
\end{tabular}
\end{table}

Il LoRA \emph{non} supera mai il baseline, a nessuna dimensione: le due curve di
AUC restano sovrapposte (Figura~\ref{fig:datasize}). Il modello base
\emph{migliora} con più dati (da $0{,}849$ a $0{,}869$), e il LoRA lo insegue
senza mai staccarlo. Alla dimensione massima ($40\,000$, circa $80$ volte la più
piccola) il LoRA anzi peggiora leggermente AUC, F1 ed ECE, con una dinamica di
training instabile (la AUC di validazione, massima alla quinta epoca, poi crolla
fino a $\approx 0{,}63$): solo l'early stopping ne preserva le prestazioni.
\textbf{Anche l'ipotesi della scarsità di dati è quindi rigettata}: più dati
giovano al modello, non all'adattamento.

\begin{figure}[h]
\centering
\begin{tikzpicture}
\begin{axis}[
    width=0.78\textwidth, height=6cm,
    xlabel={Dimensione del training set (scala log)},
    ylabel={AUC-ROC sul test fisso},
    xmode=log, log basis x=10,
    xmin=400, xmax=50000, ymin=0.845, ymax=0.872,
    xtick={500,2000,10000,40000},
    xticklabels={500,2000,10000,40000},
    ytick={0.85,0.855,0.86,0.865,0.87},
    yticklabel style={/pgf/number format/fixed, /pgf/number format/precision=3},
    grid=major, grid style={gray!25},
    tick label style={font=\footnotesize},
    label style={font=\small},
    legend pos=south east, legend cell align=left,
    every axis plot/.append style={very thick},
]
\addplot[color=gray, mark=o, mark size=2.5pt] coordinates {
    (500,0.8485) (2000,0.8565) (10000,0.8658) (40000,0.8690)
};
\addplot[color=blue, mark=square*, mark size=2.5pt] coordinates {
    (500,0.8481) (2000,0.8553) (10000,0.8657) (40000,0.8674)
};
\legend{TabPFN base, TabPFN + LoRA}
\end{axis}
\end{tikzpicture}
\caption{Esperimento~3 -- AUC-ROC sul test fisso di \texttt{gmsc} al crescere del
training set. Le due curve sono praticamente sovrapposte: il LoRA segue il
baseline senza mai superarlo, a ogni quantità di dati.}
\label{fig:datasize}
\end{figure}

\subsection{Discussione del fine-tuning LoRA}
\label{sec:lora-discussione}
I tre esperimenti convergono verso una conclusione coerente
(Tabella~\ref{tab:summary}). Il LoRA fine-tuning non migliora la capacità
discriminativa di TabPFN-2.5 su classificazione binaria tabulare. Le due
spiegazioni alternative più naturali sono state escluse sperimentalmente:
\begin{itemize}[nosep]
  \item \textbf{non è un problema di capacità}: un aumento di $4\times$ dei
  parametri allenabili non produce alcun cambiamento (Sez.~\ref{sec:capacity});
  \item \textbf{non è un problema di dati}: un aumento fino a $80\times$ dei dati
  di training non produce miglioramenti, e un addestramento prolungato degrada il
  modello (Sez.~\ref{sec:datasize});
  \item \textbf{non è un problema di dominio}: la valutazione è \emph{in-domain}
  (finanziario), quindi il mancato beneficio non è imputabile a uno shift di
  dominio (Sez.~\ref{sec:exp5}).
\end{itemize}

\begin{table}[h]
\centering
\caption{Sintesi degli esperimenti LoRA e delle ipotesi testate.}
\label{tab:summary}
\small
\begin{tabular}{lll}
\toprule
Ipotesi & Esperimento & Esito \\
\midrule
Generalizzazione in-domain & Exp.~1 (finanziario $\to$ finanziario) & LoRA $\approx$ base \\
Capacità insufficiente & Exp.~2 (ablation $r$, target) & nessun cambiamento; overfitting \\
Dati insufficienti & Exp.~3 (fino a $80\times$ dati) & nessun gain; over-training degrada \\
\bottomrule
\end{tabular}
\end{table}

La spiegazione residua, e più solida, è che \textbf{i pesi pre-addestrati di
TabPFN-2.5 sono già vicini all'ottimo} per questa classe di problemi: il suo
in-context learning estrae quasi tutto il segnale disponibile, e un adattamento a
basso rango non può aggiungervi capacità discriminativa, potendo solo
allontanarsi dal buon punto di partenza. Il fatto che il risultato valga
\emph{in-domain} lo rende particolarmente netto: non è un problema di
trasferimento tra domini diversi, ma un vero soffitto del modello. L'unico
segnale ricorrente -- per quanto modesto -- è sulla \emph{calibrazione} (ECE), il
che suggerisce che eventuali benefici vadano cercati nel post-processing piuttosto
che nell'adattamento dei pesi: è esattamente la direzione esplorata nella
Sezione~\ref{sec:pipeline}.

\paragraph{Note implementative rilevanti.} Durante lo sviluppo sono emersi alcuni
accorgimenti tecnici necessari: (i)~l'early stopping è indispensabile, poiché
l'addestramento può degradare il modello; (ii)~su dataset grandi occorre
abilitare l'\emph{activation checkpointing} per non esaurire la memoria della GPU
in fase di training; (iii)~per dataset oltre il limite ufficiale va impostato
\texttt{ignore\_pretraining\_limits=True}.

%==============================================================================
\section{Pipeline esterna: ensembling, calibrazione e soglia}
\label{sec:pipeline}
%==============================================================================

\subsection{Motivazione e approccio}
Il capitolo precedente ha stabilito che il fine-tuning a basso rango non
migliora la capacità discriminativa di TabPFN-2.5: i pesi pre-addestrati sono
già vicini all'ottimo e l'unico segnale ricorrente -- per quanto modesto -- è
sul versante della \emph{calibrazione}. Questa osservazione indirizza la
seconda parte del lavoro: invece di modificare il modello, lo si tratta come un
predittore \textbf{congelato} e si interviene unicamente sul suo
\textbf{output} (le probabilità della classe positiva). È l'approccio speculare
al LoRA: il valore aggiunto va cercato \emph{attorno} al modello, su assi che
l'adattamento dei pesi non può raggiungere.

Vengono studiati due blocchi, ciascuno mirato a un diverso asse della qualità
predittiva:
\begin{itemize}[nosep]
  \item \textbf{Ensembling} (Esp.~4): combinare più predizioni per ridurre la
  varianza della stima, con l'obiettivo di migliorare la discriminazione
  (AUC-ROC);
  \item \textbf{Calibrazione e ottimizzazione della soglia} (Esp.~5):
  post-processare le probabilità per migliorarne l'affidabilità (ECE, Brier) e
  scegliere il \emph{cutoff} decisionale ottimale per l'F1 macro.
\end{itemize}

\paragraph{Quale intervento muove quale metrica.} I due blocchi sono
complementari e non ridondanti: ciascuno agisce su un asse distinto, come
riassunto nella Tabella~\ref{tab:chimuovecosa}. Questa separazione è la chiave
di lettura dei risultati: un ensembling che riordina i campioni può spostare
l'AUC; una trasformazione monotona delle probabilità (calibrazione) ne migliora
l'affidabilità senza toccarne l'ordinamento, quindi lascia l'AUC invariata; la
sola scelta della soglia agisce esclusivamente sull'F1.

\begin{table}[h]
\centering
\caption{Asse di azione atteso di ciascun intervento. ``='' indica invarianza
attesa per costruzione (es.\ una trasformazione monotona non altera l'AUC).}
\label{tab:chimuovecosa}
\small
\begin{tabular}{lcccc}
\toprule
Intervento & AUC-ROC & F1 & Brier & ECE \\
\midrule
Ensembling (Esp.~4)            & \checkmark & (\checkmark) & --         & --         \\
Calibrazione (Esp.~5a)         & $=$        & $=$          & \checkmark & \checkmark \\
Ottimizzazione soglia (Esp.~5b)& $=$        & \checkmark   & $=$        & $=$        \\
\bottomrule
\end{tabular}
\end{table}

\subsection{Metodologia sperimentale}
\label{sec:pipeline-metodo}

\paragraph{Dataset.} Gli esperimenti adottano un benchmark di sei dataset
binari: quattro di dimensione medio-piccola (\texttt{thyroid},
\texttt{blood\_transfusion}, \texttt{credit\_g}, \texttt{adult}) e due di grandi
dimensioni -- \texttt{bank\_marketing} (OpenML~1461, $45\,211$ righe) e
\texttt{magic\_telescope} (OpenML~1120, $19\,020$ righe). La ragione dei due
dataset grandi è metodologica: le metriche di calibrazione (in particolare l'ECE,
basato su 10 bin) sono rumorose su insiemi di test piccoli; i dataset
\texttt{credit\_g} ($1\,000$) e \texttt{blood\_transfusion} ($748$) sono troppo
ridotti perché le conclusioni su calibrazione e soglia siano robuste se prese
isolatamente, mentre i dataset grandi forniscono insiemi di test ampi e stime
stabili.

\paragraph{Split a tre vie e assenza di leakage.} Per ciascun dataset si adotta
una suddivisione stratificata in tre parti: \emph{train/context}
($50\%$), su cui TabPFN esegue l'in-context learning; \emph{calibration}
($25\%$), su cui si stimano il calibratore e la soglia ottimale; e \emph{test}
($25\%$), su cui si misurano tutte le metriche. Calibratore e soglia sono
fittati \emph{esclusivamente} sul set di calibrazione, mai sul test: è lo stesso
principio anti-leakage dell'early stopping del capitolo precedente. I dataset
più grandi sono campionati fino a un massimo prefissato per contenere i tempi di
inferenza su CPU, mantenendo comunque insiemi di test ampi.

\paragraph{Ripetizioni e riproducibilità.} L'intera procedura (split, fit,
valutazione) è ripetuta su più semi indipendenti ($5$ per l'Esp.~5, $3$ per
l'Esp.~4, computazionalmente più oneroso). I risultati sono riportati come
media $\pm$ deviazione standard sulle ripetizioni: ciò quantifica esplicitamente
il rumore, rilevante soprattutto sui dataset piccoli, ed evita di trarre
conclusioni da un singolo split fortunato.

\subsection{Esperimento 4: ensembling}
\label{sec:exp-ensemble}
Si confrontano tre strategie, tutte a pesi congelati:
\begin{itemize}[nosep]
  \item \texttt{base}: una singola istanza di TabPFN (riferimento);
  \item \texttt{self\_ensemble}: media delle probabilità di $5$ istanze di
  TabPFN con semi interni diversi (diverse permutazioni/preprocessing). Misura il
  guadagno ottenibile dalla sola riduzione della varianza interna del modello,
  senza ricorrere ad altri modelli;
  \item \texttt{tabpfn+gbm}: media tra TabPFN e un \texttt{GradientBoosting}
  addestrato sullo stesso contesto (ensemble eterogeneo).
\end{itemize}

\begin{table}[h]
\centering
\caption{Esperimento 4 -- ensembling. Media (deviazione standard) su $3$ semi.
In \textbf{grassetto} la migliore AUC-ROC per dataset (metrica bersaglio).
$\uparrow$ meglio se alto, $\downarrow$ meglio se basso.}
\label{tab:ensemble}
\small
\setlength{\tabcolsep}{4pt}
\begin{tabular}{llcccc}
\toprule
Dataset & Metodo & AUC-ROC$\uparrow$ & F1$\uparrow$ & Brier$\downarrow$ & ECE$\downarrow$ \\
\midrule
\multirow{3}{*}{\texttt{credit\_g}}
 & base          & \textbf{0.793 (0.046)} & 0.644 (0.024) & 0.163 (0.014) & 0.061 (0.011) \\
 & self\_ensemble& \textbf{0.793 (0.043)} & 0.660 (0.031) & 0.163 (0.013) & 0.058 (0.003) \\
 & tabpfn+gbm    & 0.790 (0.061) & 0.670 (0.055) & 0.162 (0.020) & 0.060 (0.004) \\
\addlinespace
\multirow{3}{*}{\texttt{blood\_transfusion}}
 & base          & \textbf{0.783 (0.034)} & 0.653 (0.018) & 0.144 (0.012) & 0.052 (0.009) \\
 & self\_ensemble& 0.782 (0.035) & 0.653 (0.011) & 0.144 (0.012) & 0.056 (0.004) \\
 & tabpfn+gbm    & 0.768 (0.033) & 0.658 (0.029) & 0.150 (0.014) & 0.052 (0.023) \\
\addlinespace
\multirow{3}{*}{\texttt{thyroid}}
 & base          & 0.999 (0.001) & 0.980 (0.008) & 0.005 (0.001) & 0.005 (0.000) \\
 & self\_ensemble& 0.999 (0.001) & 0.979 (0.012) & 0.005 (0.002) & 0.006 (0.000) \\
 & tabpfn+gbm    & \textbf{0.999 (0.001)} & 0.986 (0.002) & 0.004 (0.001) & 0.004 (0.001) \\
\addlinespace
\multirow{3}{*}{\texttt{adult}}
 & base          & 0.913 (0.006) & 0.795 (0.012) & 0.098 (0.003) & 0.020 (0.005) \\
 & self\_ensemble& 0.913 (0.006) & 0.792 (0.011) & 0.098 (0.003) & 0.018 (0.006) \\
 & tabpfn+gbm    & \textbf{0.915 (0.008)} & 0.795 (0.004) & 0.096 (0.004) & 0.019 (0.005) \\
\addlinespace
\multirow{3}{*}{\texttt{bank\_marketing}}
 & base          & 0.933 (0.006) & 0.740 (0.018) & 0.063 (0.004) & 0.019 (0.004) \\
 & self\_ensemble& \textbf{0.936 (0.006)} & 0.747 (0.014) & 0.062 (0.004) & 0.013 (0.002) \\
 & tabpfn+gbm    & 0.929 (0.007) & 0.725 (0.017) & 0.065 (0.004) & 0.021 (0.002) \\
\addlinespace
\multirow{3}{*}{\texttt{magic\_telescope}}
 & base          & 0.943 (0.004) & 0.876 (0.006) & 0.082 (0.004) & 0.014 (0.005) \\
 & self\_ensemble& \textbf{0.944 (0.004)} & 0.877 (0.004) & 0.081 (0.004) & 0.016 (0.005) \\
 & tabpfn+gbm    & 0.940 (0.003) & 0.872 (0.009) & 0.085 (0.003) & 0.030 (0.005) \\
\bottomrule
\end{tabular}
\end{table}

I risultati (Tabella~\ref{tab:ensemble}) indicano che \textbf{l'ensembling non
migliora la discriminazione} di TabPFN-2.5.

\paragraph{Self-ensemble: nessun guadagno oltre il rumore.} Mediare le
probabilità di cinque istanze di TabPFN produce variazioni di AUC dell'ordine di
$\pm 0{,}003$, sistematicamente \emph{interne} alla deviazione standard fra i
semi (che va da $0{,}004$ a $0{,}046$): il più ampio, su \texttt{bank\_marketing}
($0{,}933 \to 0{,}936$), resta entro un errore standard. Il motivo è strutturale:
TabPFN esegue già internamente un ensemble su più permutazioni del contesto
(parametro \texttt{n\_estimators}), quindi la varianza residua del singolo
modello è bassa e mediare ulteriori copie non aggiunge informazione. Si osserva
al più un lieve effetto di \emph{smoothing} sull'ECE in un paio di casi
(\texttt{bank\_marketing} $0{,}019 \to 0{,}013$), comunque modesto.

\paragraph{Ensemble eterogeneo: tendenzialmente controproducente.} Mediare
TabPFN con un \texttt{GradientBoosting} \emph{abbassa} l'AUC nella maggioranza
dei dataset (\texttt{blood\_transfusion} $0{,}783 \to 0{,}768$,
\texttt{magic\_telescope} $0{,}943 \to 0{,}940$, \texttt{bank\_marketing}
$0{,}933 \to 0{,}929$): essendo il GBM mediamente più debole di TabPFN su questi
problemi, la media viene trascinata verso il modello peggiore. L'unica eccezione
è \texttt{thyroid}, dove il GBM è competitivo e il mix ottiene il miglior F1
($0{,}986$); su \texttt{adult} il vantaggio di AUC ($+0{,}002$) è marginale.
In assenza di un peso adattivo o di un modello esterno almeno pari a TabPFN,
l'ensemble eterogeneo non è una strategia affidabile.

\subsection{Esperimento 5: calibrazione e ottimizzazione della soglia}
\label{sec:exp-calib}
A partire dalle probabilità grezze di TabPFN si confrontano quattro varianti:
\begin{itemize}[nosep]
  \item \texttt{base}: probabilità grezze, soglia $0{,}5$ (riferimento);
  \item \texttt{platt}: calibrazione con \emph{Platt scaling} (regressione
  logistica a una variabile sulle probabilità di calibrazione);
  \item \texttt{isotonic}: calibrazione con \emph{regressione isotonica}
  (mappatura monotona non parametrica, più flessibile ma più avida di dati);
  \item \texttt{threshold}: probabilità grezze ma con soglia scelta per
  massimizzare l'F1 macro sul set di calibrazione.
\end{itemize}

\begin{table}[h]
\centering
\caption{Esperimento 5 -- calibrazione e soglia. Media (deviazione standard)
su $5$ semi. In \textbf{grassetto} il miglior F1 e il miglior ECE per dataset.
$\uparrow$ meglio se alto, $\downarrow$ meglio se basso.}
\label{tab:calib}
\small
\setlength{\tabcolsep}{4pt}
\begin{tabular}{llcccc}
\toprule
Dataset & Metodo & AUC-ROC$\uparrow$ & F1$\uparrow$ & Brier$\downarrow$ & ECE$\downarrow$ \\
\midrule
\multirow{4}{*}{\texttt{credit\_g}}
 & base      & 0.801 (0.034) & 0.657 (0.025) & 0.161 (0.010) & 0.061 (0.011) \\
 & platt     & 0.801 (0.034) & 0.618 (0.019) & 0.167 (0.006) & 0.084 (0.021) \\
 & isotonic  & 0.789 (0.033) & 0.681 (0.044) & 0.166 (0.016) & \textbf{0.059 (0.037)} \\
 & threshold & 0.801 (0.034) & \textbf{0.710 (0.035)} & 0.161 (0.010) & 0.061 (0.011) \\
\addlinespace
\multirow{4}{*}{\texttt{blood\_transfusion}}
 & base      & 0.743 (0.066) & 0.644 (0.019) & 0.153 (0.017) & 0.073 (0.031) \\
 & platt     & 0.743 (0.066) & 0.451 (0.033) & 0.160 (0.005) & 0.073 (0.020) \\
 & isotonic  & 0.733 (0.061) & 0.606 (0.091) & 0.156 (0.011) & \textbf{0.057 (0.024)} \\
 & threshold & 0.743 (0.066) & \textbf{0.678 (0.036)} & 0.153 (0.017) & 0.073 (0.031) \\
\addlinespace
\multirow{4}{*}{\texttt{thyroid}}
 & base      & 0.999 (0.001) & 0.977 (0.008) & 0.005 (0.001) & 0.006 (0.001) \\
 & platt     & 0.999 (0.001) & 0.977 (0.006) & 0.006 (0.001) & 0.017 (0.003) \\
 & isotonic  & 0.992 (0.008) & 0.978 (0.007) & 0.006 (0.001) & \textbf{0.005 (0.001)} \\
 & threshold & 0.999 (0.001) & \textbf{0.978 (0.008)} & 0.005 (0.001) & 0.006 (0.001) \\
\addlinespace
\multirow{4}{*}{\texttt{adult}}
 & base      & 0.917 (0.007) & 0.798 (0.011) & 0.096 (0.004) & 0.019 (0.004) \\
 & platt     & 0.917 (0.007) & 0.791 (0.011) & 0.098 (0.004) & 0.046 (0.003) \\
 & isotonic  & 0.915 (0.007) & 0.790 (0.021) & 0.097 (0.004) & \textbf{0.019 (0.009)} \\
 & threshold & 0.917 (0.007) & \textbf{0.804 (0.010)} & 0.096 (0.004) & 0.019 (0.004) \\
\addlinespace
\multirow{4}{*}{\texttt{bank\_marketing}}
 & base      & 0.931 (0.005) & 0.742 (0.015) & 0.063 (0.003) & 0.019 (0.003) \\
 & platt     & 0.931 (0.005) & 0.732 (0.019) & 0.067 (0.003) & 0.034 (0.004) \\
 & isotonic  & 0.929 (0.006) & 0.739 (0.030) & 0.064 (0.003) & \textbf{0.015 (0.006)} \\
 & threshold & 0.931 (0.005) & \textbf{0.784 (0.006)} & 0.063 (0.003) & 0.019 (0.003) \\
\addlinespace
\multirow{4}{*}{\texttt{magic\_telescope}}
 & base      & 0.944 (0.005) & 0.877 (0.004) & 0.081 (0.004) & \textbf{0.014 (0.003)} \\
 & platt     & 0.944 (0.005) & \textbf{0.878 (0.005)} & 0.083 (0.004) & 0.029 (0.006) \\
 & isotonic  & 0.943 (0.005) & 0.877 (0.005) & 0.081 (0.003) & 0.018 (0.005) \\
 & threshold & 0.944 (0.005) & 0.877 (0.004) & 0.081 (0.004) & \textbf{0.014 (0.003)} \\
\bottomrule
\end{tabular}
\end{table}

I risultati (Tabella~\ref{tab:calib}) raccontano una storia netta, in parte
diversa dalle attese.

\paragraph{La soglia è il guadagno reale, concentrato sugli sbilanciati.}
L'ottimizzazione della soglia migliora l'F1 macro su tutti e sei i dataset
lasciando \emph{esattamente} invariate AUC, Brier ed ECE -- come previsto dalla
Tabella~\ref{tab:chimuovecosa}, poiché agisce solo sulla regola di decisione e
non sulle probabilità. L'entità del guadagno scala con lo sbilanciamento: è
sostanziale su \texttt{credit\_g} ($F1: 0{,}657 \to 0{,}710$, $+5{,}3$ punti) e
\texttt{bank\_marketing} ($0{,}742 \to 0{,}784$, $+4{,}2$ punti, classe
positiva al $12\%$), apprezzabile su \texttt{blood\_transfusion}
($0{,}644 \to 0{,}678$) e \texttt{adult} ($0{,}798 \to 0{,}804$), e nullo sui
dataset bilanciati o saturi (\texttt{magic\_telescope}, \texttt{thyroid}), dove
la soglia $0{,}5$ è già ottimale. È il risultato positivo più chiaro e robusto
dell'intero progetto: un intervento a costo nullo che corregge il \emph{cutoff}
decisionale là dove l'asimmetria di classe rende l'F1 a soglia fissa
sub-ottimale.

\paragraph{La calibrazione non aiuta, perché TabPFN è già ben calibrato.}
Il dato più informativo è l'ECE del modello \texttt{base}: $0{,}019$ su
\texttt{adult} e \texttt{bank\_marketing}, $0{,}014$ su \texttt{magic\_telescope},
$0{,}006$ su \texttt{thyroid}. Sono valori già molto bassi: le probabilità di
TabPFN sono affidabili \emph{senza} alcun post-processing. Di conseguenza la
ricalibrazione ha poco margine e, anzi, tende a peggiorare: il \textbf{Platt
scaling} aumenta l'ECE quasi ovunque (es.\ \texttt{adult} $0{,}019 \to 0{,}046$,
\texttt{thyroid} $0{,}006 \to 0{,}017$) perché impone una sigmoide a
probabilità che già non ne hanno bisogno, e sui dataset piccoli ne distorce
gravemente l'F1 (\texttt{blood\_transfusion} $0{,}644 \to 0{,}451$). La
\textbf{regressione isotonica} è più prudente e ottiene il miglior ECE in
cinque dataset su sei, ma con margini minimi (es.\ \texttt{bank\_marketing}
$0{,}019 \to 0{,}015$) e al prezzo di un lieve calo di AUC, dovuto ai
\emph{gradini} della mappatura che introducono pareggi nel \emph{ranking}
(\texttt{thyroid} $0{,}999 \to 0{,}992$). Nessuno dei due giustifica
l'adozione come passo standard.

\paragraph{Sintesi.} Degli interventi di post-processing, solo
l'ottimizzazione della soglia produce un beneficio consistente, e unicamente
sull'F1 dei dataset sbilanciati. La calibrazione è superflua perché il modello
è già calibrato: un risultato che \emph{estende} la conclusione del capitolo
LoRA -- dove l'unico segnale era proprio sull'ECE -- mostrando che quel margine,
misurato qui rigorosamente, è in realtà trascurabile.

\subsection{Discussione del capitolo}
\label{sec:pipeline-discussione}
I due esperimenti convergono con il capitolo precedente verso un'unica tesi
coerente. Degli interventi esterni considerati -- che agiscono sull'output senza
toccare i pesi -- due su tre \emph{non} producono alcun miglioramento robusto:
\begin{itemize}[nosep]
  \item l'\textbf{ensembling} non migliora l'AUC, perché la varianza del
  singolo modello è già bassa (TabPFN ensembla internamente) e mediare con un
  modello esterno più debole è dannoso (Sez.~\ref{sec:exp-ensemble});
  \item la \textbf{calibrazione} non migliora l'ECE, perché le probabilità di
  TabPFN sono già affidabili: il Platt scaling le peggiora e la regressione
  isotonica le ritocca solo marginalmente (Sez.~\ref{sec:exp-calib}).
\end{itemize}
L'unico intervento con un effetto reale e prevedibile è l'\textbf{ottimizzazione
della soglia} per l'F1 macro, efficace là dove lo sbilanciamento di classe rende
sub-ottimale il \emph{cutoff} di default ($+4$--$5$ punti di F1 su
\texttt{bank\_marketing} e \texttt{credit\_g}). È significativo che si tratti
dell'unico intervento che non modifica né i pesi né le probabilità, ma soltanto
la \emph{regola di decisione} che le traduce in classi.

Letta insieme al capitolo sul LoRA, l'evidenza complessiva è netta: TabPFN-2.5
è vicino all'ottimo non solo nella capacità discriminativa, ma anche nella
\emph{calibrazione} e nella \emph{varianza} delle sue stime. Né l'adattamento dei
pesi (LoRA), né l'ensembling, né la ricalibrazione aggiungono valore; il debole
segnale sull'ECE intravisto nel fine-tuning, qui misurato con ripetizioni e
deviazioni standard, si rivela trascurabile. Il margine di miglioramento residuo
non è nel modello né nelle probabilità, ma nella decisione finale: scegliere la
soglia in funzione della metrica obiettivo e dello sbilanciamento. È una
conclusione metodologicamente solida e, sul piano pratico, operativa.

%==============================================================================
\section{Conclusioni generali}
\label{sec:conclusioni}
%==============================================================================
Il progetto ha implementato e validato un sistema completo per migliorare
TabPFN-2.5 su classificazione binaria lungo due direzioni complementari, senza
mai riaddestrarne i pesi principali: un adattamento interno tramite LoRA
fine-tuning e una pipeline esterna di post-processing. L'infrastruttura
realizzata comprende il caricamento e il preprocessing robusto dei dati, quattro
metriche di valutazione, l'iniezione e l'addestramento degli adapter LoRA con
early stopping, e gli helper per ensembling, calibrazione e ottimizzazione della
soglia.

Sul piano sperimentale, cinque esperimenti con ripetizioni e ablation
convergono verso un'unica conclusione: \textbf{TabPFN-2.5 è vicino all'ottimo non
solo nella discriminazione, ma anche nella calibrazione e nella varianza delle
sue stime}. Ne consegue che:
\begin{itemize}[nosep]
  \item il LoRA fine-tuning non migliora la discriminazione, e ciò non dipende
  né dalla capacità degli adapter né dalla quantità di dati;
  \item l'ensembling non riduce ulteriormente la varianza, già bassa per via
  dell'ensemble interno del modello;
  \item la ricalibrazione delle probabilità è superflua, perché esse sono già
  ben calibrate (ECE di base $\le 0{,}02$ sui dataset grandi).
\end{itemize}

L'unico intervento con un effetto reale e prevedibile è l'\textbf{ottimizzazione
della soglia decisionale} per l'F1 macro, che produce guadagni di $4$--$5$ punti
sui dataset sbilanciati. È un risultato dal valore sia teorico sia pratico: il
margine di miglioramento residuo per un \emph{foundation model} tabulare già
maturo non risiede nel modello né nelle sue probabilità, ma nella regola di
decisione che le traduce in classi, da scegliere in funzione della metrica
obiettivo e dello sbilanciamento. I risultati negativi sulle tecniche di
adattamento, ottenuti con metodologia rigorosa (split anti-leakage, ripetizioni
multiple, deviazioni standard), sono essi stessi un contributo informativo
sullo stato dell'arte dei modelli tabulari pre-addestrati.

%==============================================================================
\appendix
\section{Struttura del codice}
%==============================================================================
\begin{itemize}[nosep]
  \item \texttt{utils/data\_loader.py} -- caricamento, caching, encoding e split
  dei dataset OpenML.
  \item \texttt{evaluation/metrics.py} -- AUC-ROC, F1 macro, Brier, ECE; tabelle
  comparative e salvataggio CSV.
  \item \texttt{models/lora.py} -- implementazione di LoRA (\texttt{LoRALinear},
  iniezione, salvataggio/caricamento, merge).
  \item \texttt{models/tabpfn\_lora.py} -- integrazione con TabPFN: costruzione
  del classificatore con adapter, loop di training, early stopping, esperimenti
  (\texttt{run\_lora\_exp5}, \texttt{run\_lora\_ablation},
  \texttt{run\_lora\_datasize\_financial}).
  \item \texttt{pipeline/external\_pipeline.py} -- pipeline esterna: helper puri
  (\texttt{average\_probas}, \texttt{fit\_calibrator}, \texttt{optimize\_threshold})
  ed esperimenti (\texttt{run\_ensemble\_experiment},
  \texttt{run\_calibration\_threshold\_experiment}); split a tre vie e
  aggregazione media$\pm$std su più semi.
\end{itemize}

\noindent\textit{Codice e configurazione completi sono versionati nel repository
del progetto. Il modello TabPFN-2.5 è eseguito localmente; il primo download dei
pesi richiede un token di licenza Prior Labs.}

\end{document}
